\chapter{Esettanulmányok ismert sérülékenységek alapján}

A következő esettanulmányok azt vizsgálják, hogy
a seccomp és a Landlock által képviselt operációs rendszer szintű
korlátozási elvek milyen módon csökkenthették volna egy sikeres sérülékenység
hatását. A cél nem a sérülékenység kiváltó okának megszüntetése, hanem a
kihasználás után elérhető fájlrendszer-hozzáférések és rendszerhívások körének
szűkítése.

\section{CVE-2022-44268: fájlolvasás képfeldolgozás közben}

A CVE-2022-44268 az ImageMagick képfeldolgozó szoftvert érintő
információszivárgási sérülékenység. A Red Hat leírása szerint a hiba lehetővé
teszi, hogy egy támadó képfeldolgozás közben fájlokat olvasson a szerverről:
\begin{quote}
\foreignlanguage{english}{``This flaw allows an attacker to read arbitrary
files from a server when parsing an image''} \cite{redhat-cve-2022-44268}
\end{quote}

A sérülékenység PNG képek feldolgozásakor jelentkezett. Az NVD leírása alapján
az ImageMagick 7.1.0-49 verziója információszivárgásra volt sérülékeny: PNG kép
feldolgozásakor a létrejövő képbe bekerülhetett egy tetszőleges fájl tartalma,
amennyiben a \code{magick} folyamatnak volt jogosultsága az adott fájl
olvasására \cite{nvd-cve-2022-44268}. Ez különösen olyan környezetben jelent
kockázatot, ahol a szerver nem megbízható forrásból származó képeket dolgoz fel,
például feltöltött képek átméretezése vagy konvertálása során.

A hiba működésének lényege, hogy a támadó által készített bemeneti kép
feldolgozása során a sérülékeny program olyan fájl megnyitására vehető rá,
amelyet a normál alkalmazáslogika nem kívánt volna feldolgozni. A sikeres
kihasználás hatását ezért jelentősen befolyásolja, hogy maga a képfeldolgozó
folyamat milyen fájlrendszer-hozzáférésekkel rendelkezik.

Egy Landlockhoz hasonló fájlrendszer-szintű sandbox ebben az esetben
kárcsökkentő védelmi rétegként működhetne. A képfeldolgozó folyamatot még a nem
megbízható bemenet feldolgozása előtt olyan környezetben lehetne elindítani,
ahol csak az előre meghatározott bemeneti, kimeneti és ideiglenes
munkakönyvtárakhoz fér hozzá. Ebben a modellben a Landlock szabályok nem a
bemeneti kép által hivatkozott útvonalakból származnak, hanem a
futtatókörnyezet konfigurációjából. Így ha a sérülékenység hatására a program
megpróbálna például rendszerfájlokat megnyitni, a kernel a Landlock szabályok
alapján megtagadná a hozzáférést.

A seccomp ebben az esetben kiegészítő védelmi rétegként jelenhet meg. Egy
képfeldolgozó worker működéséhez általában nincs szükség például új folyamatok
indítására vagy hálózati kommunikációra. Egy megfelelő seccomp szűrő ezért
tilthatja azokat a rendszerhívásokat, amelyek nem szükségesek a képfeldolgozás
normál működéséhez. Ez nem akadályozza meg magát a CVE-2022-44268 jellegű
alkalmazáshibát, de csökkentheti a sikeres kihasználás után végrehajtható
további műveletek körét.

\section{CVE-2021-41773: path traversal webszerver környezetben}

A CVE-2021-41773 az Apache HTTP Server 2.4.49 verzióját érintő
útvonalkezelési sérülékenység. Az NVD leírása szerint a hiba az Apache HTTP
Server 2.4.49 útvonal-normalizálási változtatásához kapcsolódott, és lehetővé
tette, hogy a támadó path traversal támadással URL-eket képezzen le a várt
dokumentumgyökéren kívüli fájlokra \cite{nvd-cve-2021-41773}. Az Apache saját
biztonsági tájékoztatója szintén az Apache HTTP Server 2.4 sérülékenységei
között tartja nyilván a hibát \cite{apache-httpd-vulnerabilities-24}.

A sérülékenység lényege, hogy a webszerver hibásan kezelhetett bizonyos
útvonalkomponenseket, ezért egy speciálisan összeállított kérés a kiszolgálásra
szánt könyvtáron kívüli fájlhoz vezethetett. Az NVD leírása alapján a kérés
akkor lehetett sikeres, ha a dokumentumgyökéren kívüli fájlokat más
konfigurációs szabály, például a szokásos ``require
all denied'' beállítás nem védte. Bizonyos konfigurációkban, CGI szkriptek
engedélyezése esetén a hiba akár távoli kódfuttatáshoz is vezethetett
\cite{nvd-cve-2021-41773}.

Ez a sérülékenység jól mutatja, hogy egy webszerver esetén az alkalmazásszintű
útvonalellenőrzés hibája milyen közvetlen fájlrendszer-hozzáférési problémához
vezethet. Egy statikus fájlokat kiszolgáló folyamatnak normál esetben nincs
szüksége a teljes fájlrendszer olvasására. A működéséhez tipikusan elegendő,
ha hozzáfér a dokumentumgyökérhez, például egy \code{/var/www} alatti
könyvtárhoz, valamint esetleg néhány előre meghatározott konfigurációs vagy
ideiglenes könyvtárhoz.

A Landlock ebben az esetben nem az útvonal-normalizálási hibát javítaná ki,
hanem a hibás útvonalkezelés következményét korlátozhatná. Ha a webszerver
folyamat már a kérések feldolgozása előtt Landlock szabályokkal csak a
dokumentumgyökér olvasására lenne jogosult, akkor egy path traversal jellegű
hiba esetén a dokumentumgyökéren kívüli fájlok megnyitását a kernel
megtagadhatná. Az alkalmazás hibás logikája eljuthatna ugyan a fájlmegnyitási
műveletig, de az operációs rendszer szintű hozzáférés-ellenőrzés miatt például
a \code{/etc}, \code{/home}, vagy más nem engedélyezett könyvtárak tartalma nem
válna olvashatóvá a folyamat számára.

Ez a megközelítés különösen jól kapcsolódik a dolgozatban bemutatott
\code{http\_\allowbreak server} demonstrációs példaprogramhoz. A példaprogramban
a Landlock szabályrendszer a kiszolgálható fájlokat a \code{www/} könyvtárra
korlátozza. Ennek célja éppen az, hogy akkor is legyen egy kernel-szintű
védelmi réteg, ha a szerver alkalmazásszintű útvonalkezelése hibás lenne. A
kiszolgáló így nem a teljes fájlrendszerre kap olvasási jogot, hanem csak arra
a részre, amelyre a normál működéséhez ténylegesen szüksége van.

A seccomp ebben a példában kiegészítő szerepet tölthet be. Egy egyszerű
statikus webszervernek szüksége van hálózati kommunikációra, alapvető be- és
kimeneti műveletekre, valamint a kiszolgált fájlok megnyitására és olvasására.
Ugyanakkor nem feltétlenül szükséges számára például új folyamatok indítása,
tetszőleges programok végrehajtása vagy olyan rendszerhívások használata,
amelyek nem kapcsolódnak közvetlenül a fájlkiszolgáláshoz. Egy megfelelően
összeállított seccomp szűrő ezért tovább csökkentheti a támadó által
kihasználható műveletek körét.

\section{Az esettanulmányok tanulságai}

A bemutatott két sérülékenység közös jellemzője, hogy a sikeres kihasználás
hatása erősen függött a sérülékeny folyamat jogosultságaitól. Az ImageMagick
esetében az volt lényeges, hogy a képfeldolgozó folyamat milyen fájlokat tudott
olvasni. Az Apache HTTP Server esetében pedig az, hogy a webszerver folyamat
hozzáférhetett-e a dokumentumgyökéren kívüli útvonalakhoz.

Mindkét esetben jól látható, hogy az alkalmazásszintű hibák mögé érdemes
operációs rendszer szintű korlátozásokat is helyezni. A Landlock a
fájlrendszer-hozzáférések körét szűkítheti le előre meghatározott útvonalakra,
míg a seccomp a használható rendszerhívásokat korlátozhatja. Ezek a
mechanizmusok nem akadályozzák meg szükségszerűen a sérülékenység kiváltó okát,
de csökkenthetik azt, hogy egy sikeres kihasználás után a támadó milyen
erőforrásokat érhet el.

A legfontosabb tanulság az, hogy a sandbox szabályait nem a nem megbízható
bemenetből kell levezetni, hanem a program legitim működéséhez szükséges
erőforrásokból. Egy képfeldolgozó komponens esetén ilyen lehet a bemeneti és
kimeneti munkakönyvtár, egy statikus webszerver esetén pedig a
dokumentumgyökér. Ha a folyamat már a nem megbízható bemenet feldolgozása előtt
ilyen korlátozott környezetben fut, akkor egy alkalmazáshiba hatása kisebb
területre szorítható.

Ez a megközelítés a többrétegű védekezés elvéhez illeszkedik. A sérülékeny
szoftvert továbbra is frissíteni kell, az alkalmazásszintű hibákat pedig
javítani kell, de ezek mellett a folyamat jogosultságainak csökkentése
önállóan is értékes védelmi réteg lehet.
